\documentclass[11pt]{article}
\usepackage[margin=1.1in]{geometry}
\usepackage{amsmath,amssymb,amsthm}
\usepackage{booktabs}
\usepackage[colorlinks=true,linkcolor=blue,citecolor=blue,urlcolor=blue]{hyperref}

\newtheorem{theorem}{Theorem}
\newtheorem{lemma}{Lemma}
\newtheorem{remark}{Remark}
\newtheorem{corollary}{Corollary}

\newcommand{\R}{\mathbb{R}}
\newcommand{\Z}{\mathbb{Z}}
\newcommand{\conv}{\ast}
\newcommand{\code}[1]{\texttt{#1}}

\title{Commutation of PSF Convolution and Grism Dispersion:\\
Justification for the Convolve-After-Dispersion Fast Path}
\author{kl\_roman\_pipe -- \code{docs/derivations/psf\_dispersion\_commutation.tex}}
\date{2026-07-04 (draft for review)}

\begin{document}
\maketitle

\begin{abstract}
The grism forward model currently applies PSF convolution to every wavelength
slice of the spectral datacube before dispersing (per-eval cost: $2N_\lambda$
FFTs, measured 45\% of the grism leg at the flagship configuration). We prove
that, for a wavelength-independent PSF --- the assumption the pipeline already
makes (\code{source.py}, issue \#51) --- per-slice convolution followed by
dispersion is mathematically identical, on the infinite plane, to dispersion
followed by a \emph{single} convolution of the 2D dispersed image. On the
finite stamp the two orderings differ only by boundary-truncation terms, which
we bound explicitly by near-boundary flux of the same class already budgeted
by \code{folding\_threshold}. We give the generalization to per-line PSFs
(preserving a path to wavelength-dependent PSFs), a first-order estimate of
the error induced by assuming PSF constancy across a single emission line, and
the empirical validation protocol required before the fast path may become a
default.
\end{abstract}

\section{Setup and notation}
\label{sec:setup}

All operators below name concrete code paths at branch \code{se/speedups}
(commit \code{5ec0030} era); line numbers are indicative.

Let $W \subset \Z^2$ be the fine spatial stamp (flagship: $96\times96$, the
$32\times32$ detector stamp at spatial oversample 3). The spectral cube is a
finite family of slices $c = \{c_k\}_{k=1}^{K}$, $c_k : W \to \R$, at coarse
wavelengths $\lambda_k$ ($K = N_\lambda = 25$ at flagship), produced by
\code{build\_cube} (\code{source.py:336--474}).

\paragraph{Extension and restriction.} $E : \R^W \to \R^{\Z^2}$ denotes
extension by zero; $R : \R^{\Z^2} \to \R^W$ denotes restriction to $W$. Both
are linear.

\paragraph{PSF convolution $P$.} For a kernel $p \ge 0$ with
$\sum p = 1$, define on the infinite lattice $P f = p \conv f$. The
implemented operator (\code{psf.py:convolve\_fft}, \code{bin=False}) is
\[
\tilde P = R \circ P \circ E ,
\]
i.e.\ zero-pad to $M \ge 2N - 1$ (flagship: $192 \ge 2\cdot96-1$ via
\code{next\_fast\_len}), multiply FFTs, inverse-transform, crop. Because
$M \ge 2N-1$, the circular convolution equals the linear convolution on the
padded domain: $\tilde P$ is \emph{exactly} $R\,P\,E$, with no wrap-around
term. This exactness is load-bearing for everything below.

\paragraph{Sub-pixel shift $B_d$.} Dispersion shifts each slice by a real
displacement $d_k = (\Delta y_k, \Delta x_k)$ with
$\Delta x_k = s_k \cos\alpha$, $\Delta y_k = s_k \sin\alpha$,
$s_k = (\lambda_k - \lambda_{\rm ref})/D \cdot o$ (\code{dispersion.py:173--211};
$D$ = dispersion in nm per coarse pixel, $o$ = spatial oversample). The
implementation is pull-semantics bilinear interpolation
(\code{jax.scipy.ndimage.map\_coordinates}, \code{order=1},
\code{mode='constant'}, \code{cval=0}). Write $d = n + \delta$ with
$n \in \Z^2$, $\delta \in [0,1)^2$. On the infinite lattice, bilinear
interpolation is exactly
\[
B_d f = \tau_n\!\left( t_\delta \conv f \right),
\qquad
t_\delta = \begin{pmatrix}
(1-\delta_y)(1-\delta_x) & (1-\delta_y)\,\delta_x\\
\delta_y\,(1-\delta_x) & \delta_y\,\delta_x
\end{pmatrix},
\]
where $\tau_n$ is translation by the integer vector $n$ and $t_\delta$ is the
$2\times2$ tent (triangle) kernel with $\sum t_\delta = 1$. That is:
\emph{bilinear shift is itself a convolution followed by a translation}. The
implemented operator is $\tilde B_d = R \circ B_d \circ E$.

\paragraph{Dispersion $D$.} With throughputs $w_k = t_k\,\Delta\lambda > 0$
(\code{dispersion.py:220}),
\[
D c \;=\; \sum_{k=1}^{K} w_k\, B_{d_k} c_k ,
\qquad
\tilde D c \;=\; \sum_{k=1}^{K} w_k\, \tilde B_{d_k} c_k .
\]

\paragraph{The two orderings.} The current path (A) and proposed fast path
(B) for the pre-pixel-response dispersed image are
\[
G_{\rm A} \;=\; \sum_k w_k\, \tilde B_{d_k}\, \tilde P\, c_k
\qquad\text{vs.}\qquad
G_{\rm B} \;=\; \tilde P \sum_k w_k\, \tilde B_{d_k}\, c_k .
\]
Both are followed by the post-dispersion pixel response (BoxPixel sinc
multiply + sum-bin, \code{grism.py:16--76}), which is identical in both paths
and is itself already a post-dispersion convolution --- precedent for
commuting linear operators past $D$ in this codebase.

\section{Exact commutation on the infinite plane}

\begin{lemma}[Convolution commutes with lattice translation and with other
convolutions]
\label{lem:comm}
For any $f$, kernel $p$, integer translation $\tau_n$, and kernel $t$:
$p \conv (\tau_n f) = \tau_n (p \conv f)$ and
$p \conv (t \conv f) = t \conv (p \conv f)$.
\end{lemma}
\begin{proof}
Both are immediate from the definition of discrete convolution
$(p\conv f)(x) = \sum_y p(y) f(x-y)$: substituting $f \mapsto \tau_n f$
re-indexes the sum; associativity and commutativity of convolution give the
second identity.
\end{proof}

\begin{theorem}[Commutation, infinite domain]
\label{thm:main}
For any kernel $p$, displacements $\{d_k\}$, weights $\{w_k\}$, and slices
$\{c_k\}$ on $\Z^2$,
\[
\sum_k w_k\, B_{d_k} (P c_k) \;=\; P\!\left(\sum_k w_k\, B_{d_k} c_k\right).
\]
\end{theorem}
\begin{proof}
Fix $k$ and write $d_k = n_k + \delta_k$. By the bilinear-shift decomposition
and Lemma~\ref{lem:comm},
\[
B_{d_k}(P c_k)
= \tau_{n_k}\big(t_{\delta_k} \conv (p \conv c_k)\big)
= \tau_{n_k}\big(p \conv (t_{\delta_k} \conv c_k)\big)
= p \conv \tau_{n_k}\big(t_{\delta_k} \conv c_k\big)
= P(B_{d_k} c_k).
\]
The claim follows by linearity of $P$ over the finite weighted sum.
\end{proof}

Theorem~\ref{thm:main} covers the exact discrete operations the code performs
(bilinear interpolation, arbitrary real $d_k$, arbitrary dispersion angle,
per-slice throughput weights); no continuum idealization is invoked. The
\emph{only} gap between Theorem~\ref{thm:main} and the implementation is the
finite stamp, treated next.

\section{Finite-domain error bound}
\label{sec:finite}

Define the \emph{ideal} dispersed image, in which all intermediates live on
the infinite lattice and only the final result is cropped:
\[
G_\infty \;=\; R \sum_k w_k\, B_{d_k}\, P\, (E c_k)
\;\overset{\text{Thm.~\ref{thm:main}}}{=}\;
R\, P \sum_k w_k\, B_{d_k} (E c_k).
\]

\begin{lemma}[Truncation terms]
\label{lem:trunc}
For each slice $k$, with $\mathbb{1}_{W^c}$ the indicator of the stamp
complement and $\|\cdot\|_1$ the absolute sum over all pixels,
\begin{align}
\big\| w_k \tilde B_{d_k} \tilde P c_k - w_k R\,B_{d_k} P E c_k \big\|_1
&\;\le\; w_k \,\big\| (P E c_k)\,\mathbb{1}_{W^c} \big\|_1
\;=:\; w_k\, \mu_k^{\rm A},
\label{eq:muA}\\
\big\| w_k \tilde P' \tilde B_{d_k} c_k - w_k \tilde P' R\, B_{d_k} E c_k \big\|_1
&\;=\; 0,
\qquad
\big\| w_k R P E (R B_{d_k} E c_k) - w_k R P B_{d_k} E c_k \big\|_1
\;\le\; w_k\, \mu_k^{\rm B},
\label{eq:muB}
\end{align}
where $\mu_k^{\rm B} := \| (B_{d_k} E c_k)\,\mathbb{1}_{W^c} \|_1$ and
$\tilde P'$ denotes the implemented convolution applied to the (already
restricted) dispersed sum. Consequently
\[
\| G_{\rm A} - G_\infty \|_1 \le \sum_k w_k\, \mu_k^{\rm A},
\qquad
\| G_{\rm B} - G_\infty \|_1 \le \sum_k w_k\, \mu_k^{\rm B},
\qquad
\| G_{\rm A} - G_{\rm B} \|_1 \le \sum_k w_k (\mu_k^{\rm A} + \mu_k^{\rm B}).
\]
\end{lemma}
\begin{proof}
In path A the implemented pipeline computes
$\tilde B_{d} \tilde P c_k = R\,B_{d}\,E\,R\,(P E c_k)$. The composition
$E R$ zeroes exactly the content of $P E c_k$ outside $W$; since $B_d$ and
$R$ are $\ell_1$-contractions on nonnegative decompositions, the discrepancy
from $R B_d P E c_k$ is at most the mass so zeroed, giving \eqref{eq:muA}.
In path B the analogous insertion of $E R$ acts on $B_{d} E c_k$, giving
$\mu_k^{\rm B}$; the outer $\tilde P = R P E$ applied to the restricted
dispersed image is exact linear convolution of its input by
Section~\ref{sec:setup}, contributing no additional wrap term. Summation and
the triangle inequality complete the proof.
\end{proof}

\begin{remark}[What the bound means physically]
$\mu_k^{\rm A}$ is the PSF-blurred slice flux lying outside the stamp;
$\mu_k^{\rm B}$ is the shifted (un-blurred) slice flux outside the stamp.
Both are ``flux within about one PSF support radius of the stamp boundary,
after the slice's dispersion shift.'' Slices at the wavelength-window
extremes have the largest $|d_k|$ (flagship: up to $\pm12$ coarse pixels on a
32-pixel stamp) and hence dominate the bound; those same slices carry only
line-wing and continuum flux, which is why the empirical discrepancy is far
smaller than a worst-case reading of the bound. Note both paths truncate
\emph{something}; neither is ``exact'' relative to $G_\infty$. Path B inserts
one fewer $E R$ truncation between blur and shift, so there is no a priori
reason to regard path A as the more faithful of the two.
\end{remark}

\begin{corollary}[Tolerance policy]
\label{cor:tol}
\code{RenderConfig} sizes stamps and grids so that profile flux beyond the
stamp is budgeted by \code{folding\_threshold} $f_t$ (default
$5\times10^{-3}$); the near-boundary flux entering
Lemma~\ref{lem:trunc} is the same class of quantity evaluated per shifted
slice. The A/B equivalence test therefore asserts
\[
\| G_{\rm A} - G_{\rm B} \|_1 \;\le\; C\, f_t\, F_{\rm tot},
\]
with $F_{\rm tot}$ the total dispersed flux and $C$ an order-unity constant
frozen at implementation time from the measured constant across the
validation grid (Section~\ref{sec:protocol}). Tying the tolerance to $f_t$
(rather than an ad-hoc absolute) keeps the fast-path error in lockstep with
the accuracy regime the user has already selected.
\end{corollary}

\section{Wavelength dependence of the PSF}
\label{sec:lambda}

\subsection{Status quo and exact per-line generalization}
The pipeline currently applies a \emph{single} shared \code{psf\_data} to all
$K$ slices (\code{source.py:219--223}; wavelength-dependent PSFs are deferred
--- issue \#51). Under that existing assumption, Theorem~\ref{thm:main} is an
identity: the fast path changes numerics only through
Lemma~\ref{lem:trunc}-class boundary terms, and introduces \emph{no new
physical approximation whatsoever}.

If per-line PSFs $p^{(\ell)}$ are later adopted, the fast path generalizes
exactly: build per-line sub-cubes $c^{(\ell)}$ (the line decomposition
already exists inside \code{build\_cube}'s loop over lines), disperse each,
convolve each dispersed image with its own $p^{(\ell)}$, and sum:
\[
\sum_\ell P^{(\ell)}\!\left( \sum_k w_k B_{d_k} c^{(\ell)}_k \right),
\]
i.e.\ $O(N_{\rm lines})$ convolutions instead of $O(N_\lambda)$. A fully
per-slice $p_k$ falls back to the current general path --- which is
\textbf{retained} in the code by decision (2026-07-04): the fast path is an
opt-in configuration, the per-slice path remains the general reference
implementation.

\subsection{Error of assuming PSF constancy across one line}
\label{sec:lambdaerr}

For a near--diffraction-limited PSF the size scales linearly with wavelength,
so the second moment obeys
$\delta R^2_{\rm psf} / R^2_{\rm psf} \simeq 2\,\delta\lambda/\lambda$.
An unmodeled fractional PSF second-moment error $\epsilon$ propagates to a
multiplicative shear-type bias of order
$m \sim \epsilon \cdot R^2_{\rm psf}/R^2_{\rm gal}$
(standard PSF-dilution scaling).

This was measured directly (2026-07-04) with monochromatic
\code{galsim.roman.getPSF} adaptive moments (SCA 1, cross-checked SCA 9/18
and \code{pupil\_bin} 1 vs 4; script and log in
\code{experiments/sweverett/production\_speedups/psf\_lambda\_variation/}).
Measured $\sigma_{\rm psf}(1313\,{\rm nm}) = 0.0457''$; the fractional size
scaling matches naive $\delta\lambda/\lambda$ diffraction scaling to
$<1\%$, monotonic, with no aberration surprises; ellipticity changes are
$|\delta e_{1,2}| \le 5\times10^{-5}$ everywhere. Implied $m$ below uses
$r_{\rm gal} = 0.3''$ (conservative; $0.5''$ is $\sim3\times$ smaller):

\begin{center}
\begin{tabular}{lccc}
\toprule
window & $\delta\lambda$ & measured $|\delta\sigma/\sigma|$ & implied $m$ \\
\midrule
single velocity-broadened line & $\pm1$\,nm & $7.6\times10^{-4}$ & $3.5\times10^{-5}$ \\
H$\alpha$+[NII] blend span & $\pm2.5$\,nm & $1.9\times10^{-3}$ & $8.8\times10^{-5}$ \\
full modeled cube window & $\pm13.5$\,nm & $1.02\times10^{-2}$ & $4.7\times10^{-4}$ \\
\bottomrule
\end{tabular}
\end{center}

Three observations. (i) Per-line constancy is safe by a wide margin:
$m \sim 3\times10^{-5}$ sits far below the paper's ``reasonable
performance'' bar and below the $0.1\sigma$ parameter-shift criterion for
any realistic per-galaxy posterior. (ii) Constancy across the \emph{full}
cube window (relevant to the continuum, which spans all slices) is the
dominant term at $m \sim 5\times10^{-4}$ --- still sub-$10^{-3}$, and it
bounds the pre-existing issue-\#51 approximation, not anything introduced by
the fast path; it sharpens the case for the per-line generalization when
requirement-level bias control is eventually attempted, and deserves a
proper WL-formalism treatment (beyond this order-of-magnitude scaling) if
that window ever becomes the operative one. (iii) For the paper's
self-consistent mocks (fit model $\equiv$ truth model, same PSF assumption
on both sides) the bias is identically zero by construction; the numbers
above bound model-adequacy for real data.

\section{Implementation corollary: the fast path is nearly free}
\label{sec:impl}

The post-dispersion pixel response
(\code{grism.py:\_apply\_post\_dispersion\_pixel\_response}) already performs
one \code{fft2}/\code{ifft2} pair on the dispersed fine image to apply the
BoxPixel sinc. Since the commuted PSF acts at the same point in the chain,
its transfer function can be \emph{multiplied into the same k-space pass}:
\[
\widehat{G}_{\rm out} = \widehat{G}_{\rm B,raw} \cdot \hat p \cdot
\widehat{\rm sinc},
\]
so the fast path eliminates all $2N_\lambda$ per-slice FFTs (flagship: 50)
while adding zero new transforms. Caveat for implementation: the existing
pixel-response pass operates on the $96\times96$ dispersed image without
$2N-1$ padding (sinc is applied circularly by construction of the
precomputed \code{pixel\_response\_fft}); fusing the PSF there requires
either (i) demonstrating the PSF's circular-vs-linear difference is within
the Corollary~\ref{cor:tol} budget on the padded-enough stamp, or (ii)
padding that single pass to $M \ge 2N-1$ (two $192^2$ FFTs --- still a
$25\times$ reduction). The equivalence test must be run against whichever
variant is implemented.

\section{Empirical validation protocol}
\label{sec:protocol}

Required before the fast path can be enabled anywhere (and re-run before any
default flip, which separately requires user sign-off):

\begin{enumerate}
\item \textbf{Image-level A/B equivalence.} Grid over: flagship canonical
parameters; high-shift stress case (line centered near the wavelength-window
edge, maximal $|d_k|$); high velocity gradient (large $v_{\rm circ}$, low
$\cos i$); dispersion angles $\{0, 30^\circ, 90^\circ\}$; PSF FWHM
$\{0.11'', 0.18'', 0.3''\}$; with and without continuum. Assert
$\|G_{\rm A}-G_{\rm B}\|_1 / F_{\rm tot} \le C f_t$ and report the measured
$C$ distribution; freeze $C$ in the test.
\item \textbf{Gradient-level equivalence.} Same assertion for
$\nabla_\theta$ of a scalar reduction through both paths (the fast path must
be adjoint-consistent, not just primal-consistent).
\item \textbf{Posterior-level A/B.} Seeded flagship inference with path A
vs.\ path B: shear posterior mean shifts $< 0.1\sigma$, widths within a few
percent, divergence/R-hat/ESS not degraded.
\item \textbf{Pathway preservation.} Regression test asserting the per-slice
general path remains selectable and bit-identical to its pre-refactor
behavior.
\end{enumerate}

\paragraph{Decision record.} User (2026-07-04): PSF is physically
wavelength-dependent; constancy is assumed per line only, and the per-slice
pathway must be retained for future wavelength-varying PSFs. This document
plus the protocol above constitute the required justification package.

\end{document}
